% !TEX TS-program = pdflatex
% !TEX encoding = UTF-8 Unicode

% This file is a template using the "beamer" package to create slides for a talk or presentation
% - Giving a talk on some subject.
% - The talk is between 15min and 45min long.
% - Style is ornate.

% MODIFIED by Jonathan Kew, 2008-07-06
% The header comments and encoding in this file were modified for inclusion with TeXworks.
% The content is otherwise unchanged from the original distributed with the beamer package.

\documentclass{beamer}
\setbeamercovered{invisible}



% Copyright 2004 by Till Tantau <tantau@users.sourceforge.net>.
%
% In principle, this file can be redistributed and/or modified under
% the terms of the GNU Public License, version 2.
%
% However, this file is supposed to be a template to be modified
% for your own needs. For this reason, if you use this file as a
% template and not specifically distribute it as part of a another
% package/program, I grant the extra permission to freely copy and
% modify this file as you see fit and even to delete this copyright
% notice. 


\mode<presentation>
{
  \usetheme{Warsaw}
  % or ...

  % or whatever (possibly just delete it)
}


\usepackage[english]{babel}
% or whatever

\usepackage[utf8]{inputenc}
\usepackage{qrcode}
\usepackage{hyperref}
\usepackage{ytableau}
\usepackage{longtable}
% or whatever

\usepackage{times}
\usepackage[T1]{fontenc}
\usepackage{tikz}
% Or whatever. Note that the encoding and the font should match. If T1
% does not look nice, try deleting the line with the fontenc.

%%% MATH RELATED

\title{MATH 180 Strategies of Problem Solving} 
\subtitle
{} % (optional)

\author[W.R. Casper] % (optional, use only with lots of authors)
{W.R. Casper}
% - Use the \inst{?} command only if the authors have different
%   affiliation.

\institute[California State University Fullerton] % (optional, but mostly needed)
{
  Department of Mathematics\\
  California State University Fullerton}
% - Use the \inst command only if there are several affiliations.
% - Keep it simple, no one is interested in your street address.

\subject{Talks}
% This is only inserted into the PDF information catalog. Can be left
% out. 



% If you have a file called "university-logo-filename.xxx", where xxx
% is a graphic format that can be processed by latex or pdflatex,
% resp., then you can add a logo as follows:

% \pgfdeclareimage[height=0.5cm]{university-logo}{university-logo-filename}
% \logo{\pgfuseimage{university-logo}}



% Delete this, if you do not want the table of contents to pop up at
% the beginning of each subsection:
\AtBeginSubsection[]
{
  \begin{frame}<beamer>{Outline}
    \tableofcontents[currentsection,currentsubsection]
  \end{frame}
}


% If you wish to uncover everything in a step-wise fashion, uncomment
% the following command: 

%\beamerdefaultoverlayspecification{<+->}


\begin{document}

\begin{frame}
  \titlepage
\end{frame}

\begin{frame}{Outline}
  \tableofcontents
  % You might wish to add the option [pausesections]
\end{frame}

% Since this a solution template for a generic talk, very little can
% be said about how it should be structured. However, the talk length
% of between 15min and 45min and the theme suggest that you stick to
% the following rules:  

% - Exactly two or three sections (other than the summary).
% - At *most* three subsections per section.
% - Talk about 30s to 2min per frame. So there should be between about
%   15 and 30 frames, all told.

\section{SoPS Lecture 13}

\begin{frame}{Review}
  \begin{block}{Last time}
    \begin{itemize}
\pause
      \item a story about Gauss
\pause
      \item problem solving strategies
\pause
      \item math competition problems
    \end{itemize}
  \end{block}
\end{frame}

\subsection{Connect 4}

\begin{frame}{Play the game}
  \begin{block}{Activity}
    \begin{itemize}
\pause
      \item remind yourself how to play
\pause
      \item play a some games of Connect 4 with a friend
\pause
      \item best strategy to win?
    \end{itemize}
  \end{block}
\end{frame}

\begin{frame}{Connect 4 Strategy}
  \begin{block}{Discussion}
	Who can force a win?
  \end{block}
\pause
  \begin{block}{Discussion}
	What do you think is the best first move?  What is
the best second move?  What would a complete winning strategy description look like?
  \end{block}
\end{frame}

\begin{frame}{Weakly and strongly solved games}
  \begin{block}{Definition}
    A game is
    \begin{itemize}
\pause
      \item \textbf{very weakly solved} if we know which player can force a win at the start of the game
\pause
      \item \textbf{weakly solved} if we know \emph{how} the player can force a win starting from the start of the game
\pause
      \item \textbf{strongly solved} if we know which player will win and how from any possible state of the game
    \end{itemize}
  \end{block}
\end{frame}

\begin{frame}{Examples}
\pause
  \begin{block}{Example 1}
    Chomp is very weakly solved.
  \end{block}
\pause
  \begin{block}{Example 2}
    Checkers is weakly solved.
  \end{block}
\pause
  \begin{block}{Example 3}
    Connect 4 is strongly solved (89.6 GB data).
  \end{block}
\end{frame}

\begin{frame}{Weak solution}
  See weak solution on \href{https://www.youtube.com/watch?v=KaljD3Q3ct0}{YouTube}.\\
  Play against weak solution on \href{https://2swap.github.io/WeakC4/?pos=4}{Github}.
\end{frame}

\subsection{Wrapping up}

\begin{frame}{Final thoughts}
  \begin{block}{Reflection}
    \begin{itemize}
\pause
      \item What is the difference between a weak solution and a strong solution of a game?
\pause
      \item How was a graph used to visualize a weak solution of Connect 4?
    \end{itemize}
  \end{block}
\begin{center}
Submit your responses by \textbf{Tonight!}
\end{center}
\end{frame}


\end{document}


